% =====================================================================
%  第 9 章：Symanzik 改进与重整化群作用量
%  Gattringer & Lang, "Quantum Chromodynamics on the Lattice" (LNP 788)
%  Reconstructed from extract/ch09.txt; equations auto-numbered.
% =====================================================================

\chapter{Symanzik 改进与重整化群作用量}

在格点上引入 QCD 作用量时，我们必须对连续作用量中出现的导数项进行离散化。当时已经指出，任何离散化，例如费米子作用量中一阶导数的对称差分，都会产生离散化效应。通常，费米子的离散化效应为 $\order{a}$，规范场则为 $\order{a^2}$。这些效应只有在格距 $a$ 趋于零的连续极限中才会消失。然而，实现连续极限并非易事。随着 $a$ 减小，格点数目必须增加，以使物理体积保持不变（理想情况下，应当先让格点数目趋于无穷，再让 $a$ 趋于零）。因此在数值模拟中，人们总是在有限的 $a$ 下工作，必须处理离散化误差，例如将其纳入向 $a$ 趋于零的外推之中。

处理该问题的一种优雅方式是系统地减小离散化误差。我们已提到，我们所选的离散化并不是唯一的。其他离散化方式也给出相同的形式连续极限。特别是，可以把不同的项组合起来，得到离散化效应减小的格点作用量。例如，向 Wilson 费米子作用量加上一个额外项，并适当调节其系数，就可以把离散化误差从 $\order{a}$ 减小到 $\order{a^2}$。类似地，也有可能（并且对于完全的改进而言是必要的）减小所用可观测量的离散化误差。这些思想的系统化实施就是 Symanzik 改进纲领，我们将在第 9.1 节中讨论。

一个概念上不同的方法使用由重整化群（RG）变换构造的作用量。改进是通过积掉（阻塞）短距离自由度并把它们的影响包含在阻塞场的作用量中来实现的。这一策略将在第 9.2--9.4 节中介绍。在讨论了格点作用量可解析计算的自由费米子的阻塞方法（第 9.2 节）之后，我们在第 9.3 节中介绍完整 QCD 的普遍框架。最后，在第 9.4 节中我们讨论 RG 方法的一个重要应用，即识别连续对称性在格点上的对应物，并作为一个应用推导出 Ginsparg--Wilson 方程。

\section{Symanzik 改进纲领}
\label{sec:9.1}

\subsection{一个玩具例子}
\label{sec:9.1.1}

让我们用一个玩具例子开始关于改进的讨论，该例子已经包含了改进格点 QCD 时将要采取的大部分步骤。我们考虑单实变量 $x$ 的某个函数 $f(x)$ 的导数 $f'(x)$ 的（对称）离散化：
\begin{equation}
  \frac{f(x+a) - f(x-a)}{2a}
  = f'(x) + a^2\,C^{(2)}(x) + a^4\,C^{(4)}(x) + O(a^6) .
\end{equation}
注意，由于左端是（反）对称的离散化，右端只能出现 $a$ 的偶数次幂和 $f$ 的奇阶导数。我们引入了缩写 $C^{(2)}$ 和 $C^{(4)}$ 来表示修正项。由于 $a$ 具有长度量纲，为使两端量纲相同，修正项 $C^{(k)}$ 必须具有量纲 $1+k$，即 $C^{(k)} \sim \mathrm{length}^{-(1+k)}$（假定 $f$ 无量纲）。

利用泰勒级数展开
\begin{equation}
  f(x \pm a) = f(x) \pm a\,f'(x) + \frac{a^2}{2}\,f''(x)
    \pm \frac{a^3}{6}\,f'''(x) + O(a^4) ,
\end{equation}
我们可以确定首阶修正项为
\begin{equation}
  C^{(2)}(x) = \frac{1}{6}\,f'''(x) .
\end{equation}
我们在此强调，修正项的这一表达式是用连续语言给出的，即作为 $f$ 的更高阶导数。

改进的策略是，在 (9.1) 的左端加上一个离散化表达式（即由 $f(x)$、$f(x\pm a)$、$f(x\pm 2a)$、$\ldots$ 构成的表达式），使得右端的修正项在所需阶数以内被抵消。对于 $\order{a^2}$ 改进，我们于是使用如下拟设
\begin{equation}
  \frac{f(x+a) - f(x-a)}{2a} + c\,a^2\,D^{(3)}[f](x) = f'(x) + O(a^4) ,
\end{equation}
其中 $D^{(3)}[f]$ 是满足 $D^{(3)}[f] \approx f''' + O(a^2)$ 的离散化表达式，$c$ 是某个常数。

再次利用幂级数 (9.2)，容易看出
\begin{equation}
  D^{(3)}[f](x) = \frac{f(x+2a) - 2f(x+a) + 2f(x-a) - f(x-2a)}{2a^3} ,
  \qquad c = -\frac{1}{6}
\end{equation}
即可胜任，从而实现 $\order{a^2}$ 改进。然而我们指出，选择 (9.5) 并不是唯一的，例如也可以使用包含 $f(x\pm 3a)$ 的项。

让我们总结一下玩具例子中所采取的步骤，它们已经勾勒出改进格点 QCD 的方法：
\begin{itemize}
  \item 我们从所关心量的简单离散化表达式出发（在我们的例子中是导数 $f'$）。
  \item 用连续语言识别修正项（玩具例子中是更高阶导数）。
  \item 修正项具有某些对称性（例子中只有奇阶导数），并按量纲排序。
  \item 为实现改进，以适当的系数加上修正项的离散化形式，使得直到所需阶数的修正消失。
  \item 离散化修正项的选择存在一定的任意性。
\end{itemize}

完全相同的步骤和特征也出现在格点 QCD 的改进中。主要差别在于系数的确定。在玩具例子中，系数 $c$ 来自简单的代数考虑。由于 QCD 的非线性性质以及必要的重整化，QCD 中相应系数的确定要复杂得多，必须通过微扰论或非微扰匹配过程来完成。这里概述的改进方法被称为 Symanzik 改进纲领 [1--4]。

\subsection{改进格点 QCD 的框架}
\label{sec:9.1.2}

在改进格点 QCD 时，我们需要改进欧几里得关联函数，例如位于两个不同时空点 $n, m \in \Lambda$ 的插值算符 $O_1$ 和 $O_2$ 的两点函数，
\begin{equation}
  \ev{O_1(n)\,O_2(m)}
  = \frac{1}{Z}\int D[U]\,D[\psi,\bar\psi]\,
    e^{-S[U,\psi,\bar\psi]}\,O_1[U,\psi,\bar\psi;n]\,O_2[U,\psi,\bar\psi;m] ,
\end{equation}
这些关联函数可用于计算能谱和某些矩阵元。为了普遍地获得所有这些可观测量的改进，我们一般必须同时改进格点作用量 $S$ 和插值算符 $O_1$、$O_2$。

让我们从作用量的改进开始讨论。具体而言，我们在格点上以 Wilson 规范作用量 (2.49) 和带有 Wilson 狄拉克算符 (5.51) 的费米子作用量为起点。根据前面章节的讨论，我们预期作用量的费米子部分有 $\order{a}$ 的离散化误差，规范部分有 $\order{a^2}$ 的离散化误差。

按照玩具例子所示的策略，我们首先识别修正项的连续表达式。这些修正项应按量纲排序，并具有 QCD 作用量的对称性。换言之，我们写下一个有效作用量，它描述 Wilson 形式的格点 QCD 在有限 $a$ 下的行为及其趋向连续极限的过程。按照 [1, 2, 5--7]，我们把有效作用量写成如下形式
\begin{equation}
  S_{\mathrm{eff}} = \int d^4x\,
    \left[\mathcal{L}^{(0)}(x) + a\,\mathcal{L}^{(1)}(x)
    + a^2\,\mathcal{L}^{(2)}(x) + \ldots\right] .
\end{equation}
这里 $\mathcal{L}^{(0)}$ 是 (2.3) 和 (2.17) 中定义的通常的 QCD 拉格朗日量。项 $\mathcal{L}^{(k)}$（$k \geq 1$）是额外的修正项，它们由夸克场和胶子场的乘积构成，使其具有量纲 $4+k$，即 $\mathcal{L}^{(k)} \sim \mathrm{length}^{-(4+k)}$。因此，与量纲为 4 的 $\mathcal{L}^{(0)}$ 相比，这些项将包含额外的导数或夸克质量 $m$ 的幂次。要求具有格点作用量的对称性（比较第 5.4 节），可以证明首阶修正项 $\mathcal{L}^{(1)}(x)$ 可以写成下列量纲为 5 的算符的线性组合（$\sigma_{\mu\nu} \equiv [\gamma_\mu,\gamma_\nu]/2i$）：
\begin{align}
  \mathcal{L}^{(1)}_1(x)
  &= \bar\psi(x)\,\sigma_{\mu\nu}F_{\mu\nu}(x)\,\psi(x) , \nonumber\\
  \mathcal{L}^{(1)}_2(x)
  &= \bar\psi(x)\overrightarrow{D}_\mu(x)\overrightarrow{D}_\mu(x)\psi(x)
   + \bar\psi(x)\overleftarrow{D}_\mu(x)\overleftarrow{D}_\mu(x)\psi(x) , \nonumber\\
  \mathcal{L}^{(1)}_3(x)
  &= m\,\tr[F_{\mu\nu}(x)F_{\mu\nu}(x)] , \nonumber\\
  \mathcal{L}^{(1)}_4(x)
  &= m\left[\bar\psi(x)\gamma_\mu\overrightarrow{D}_\mu(x)\psi(x)
   - \bar\psi(x)\gamma_\mu\overleftarrow{D}_\mu(x)\psi(x)\right] , \nonumber\\
  \mathcal{L}^{(1)}_5(x)
  &= m^2\,\bar\psi(x)\psi(x) .
\end{align}
利用场方程 $(\gamma_\mu D_\mu + m)\psi = 0$，这一算符清单可以进一步约化，得到如下两个关系式
\begin{gather}
  \mathcal{L}^{(1)}_1 - \mathcal{L}^{(1)}_2 + 2\mathcal{L}^{(1)}_5 = 0 ,
  \nonumber\\
  \mathcal{L}^{(1)}_4 + 2\mathcal{L}^{(1)}_5 = 0 .
\end{gather}
这些关系式可用于从张成首阶修正的算符集合中消去 $\mathcal{L}^{(1)}_2$ 和 $\mathcal{L}^{(1)}_4$ 两项，并且可以证明 [6, 7]，在树图级之外（场方程成立之处）这一结论也成立。因此，只使用 $\mathcal{L}^{(1)}_1$、$\mathcal{L}^{(1)}_3$ 和 $\mathcal{L}^{(1)}_5$ 这些项就够了。其中两项 $\mathcal{L}^{(1)}_3$ 和 $\mathcal{L}^{(1)}_5$ 在相差某个因子的意义下已经存在于原来的作用量中，因此可以通过重新定义裸参数 $m$ 和 $g$（或 $\beta = 6/g^2$）来考虑它们。

因此，对于 Wilson 格点作用量的 $\order{a}$ 改进，只需加上泡利项 $\mathcal{L}^{(1)}_1$，从而得到改进作用量
\begin{equation}
  S_I = S_{\mathrm{Wilson}} + c_{sw}\,a^5\,\frac{1}{2}\sum_{\mu<\nu}\sum_{n\in\Lambda}
    \bar\psi(n)\,\sigma_{\mu\nu}\widetilde{F}_{\mu\nu}(n)\,\psi(n) .
\end{equation}
实系数 $c_{sw}$ 通常被称为 Sheikholeslami--Wohlert 系数，因为文献 [5] 的作者首次写下了改进作用量 (9.10)。如玩具例子中所述，$\widetilde{F}_{\mu\nu}$ 必须是 $\mathcal{L}^{(1)}_1$ 中相应项的离散化版本，即场强张量的格点形式。一个方便但不唯一的选择是

\begin{figure}[htbp]
  \centering
  \includegraphics[width=0.7\textwidth]{fig91.png}
  \caption*{Fig. 9.1. \emph{用于 (9.11) 中场强离散化的 $\mu$--$\nu$ 平面内基本方格之和
    $Q_{\mu\nu}(n)$ 的图形表示}.}
  \label{fig:91}
\end{figure}

\begin{equation}
  \widetilde{F}_{\mu\nu}(n) = \frac{-i}{8a^2}\left(Q_{\mu\nu}(n) - Q_{\nu\mu}(n)\right) ,
\end{equation}
其中 $Q_{\mu\nu}(n)$ 是如图 9.1 所示的 $\mu$--$\nu$ 平面内基本方格 $U_{\mu,\nu}(n)$（比较 (2.48)）之和，
\begin{equation}
  Q_{\mu\nu}(n) \equiv U_{\mu,\nu}(n) + U_{\nu,-\mu}(n) + U_{-\mu,-\nu}(n) + U_{-\nu,\mu}(n) ,
\end{equation}
它是连续场强张量的离散化。由于各项的形状使人联想到三叶草，通常把 (9.10) 中的最后一项称为三叶草项或三叶草改进。

在此我们需要补充几点说明：我们看到，对于 $\order{a}$ 改进，只需向费米子作用量添加一项就够了。唯一的纯胶子项 $\mathcal{L}^{(1)}_3$ 被发现正比于原来的规范作用量，因此被吸收进裸规范耦合的重新定义之中。相关的纯胶子算符只在量纲为 6 时才出现，即它们在 $\order{a^2}$ 阶才有贡献。这与前面关于玻色子场离散化误差仅为 $\order{a^2}$ 的论述一致。对于规范作用量的改进，我们请读者参阅原始文献，例如 [3, 8]，其中给出了所谓的 L\"uscher--Weisz 规范作用量。

第二点说明涉及满足 Ginsparg--Wilson 方程 (7.29) 的手征对称狄拉克算符的离散化误差。如果加上 $\order{a}$ 改进项，即三叶草项，所得的狄拉克算符将不再满足 Ginsparg--Wilson 方程。因此我们必须得出结论：Ginsparg--Wilson 狄拉克算符已经是 $\order{a}$ 改进的。事实上，当朴素格点狄拉克算符（它必然是任何格点狄拉克算符的组成部分）被平方时，Ginsparg--Wilson 方程 (7.29) 的非线性右端会生成泡利项 $\mathcal{L}^{(1)}_1$ 的格点离散化。$\order{a}$ 改进与手征对称性之间的这种联系已经暗示了一种确定 Sheikholeslami--Wohlert 系数的非微扰策略 [6]，我们将在讨论流的改进之后给出这一策略。

\subsection{插值算符的改进}
\label{sec:9.1.3}

我们在上一节讨论的作用量改进，对于强子质量等在壳量的 $\order{a}$ 改进已经足够了。然而，为了完整地实现关联函数 (9.6) 的 $\order{a}$ 改进（这对改进由它计算出的强子矩阵元是必要的），插值算符 $O_i$ 也需要改进。

与作用量一样，策略是找出按量纲组织、并通过要求待改进算符 $O_i$ 的对称性来识别的修正项的连续表达式。由于可以构造具有任意量子数的流，因此不可能给出一般 $n$ 点函数 $\order{a}$ 改进所需项的详尽清单。按照 [6, 7]，我们代之以讨论两个例子，即同位旋矢量轴矢流 $A_\mu^a$ 和赝标量密度 $P^a$ 的改进：
\begin{align}
  A^a_\mu(n) &= \frac{1}{2}\,\bar\psi(n)\gamma_\mu\gamma_5\tau^a\psi(n) , \nonumber\\
  P^a(n) &= \frac{1}{2}\,\bar\psi(n)\gamma_5\tau^a\psi(n) .
\end{align}
这里 $\tau^a$ 表示作用在 $\Nf = 2$ 味空间中的泡利矩阵 $\tau^1$、$\tau^2$、$\tau^3$ 之一。流 $A^a_\mu$ 的 $\order{a}$ 改进所需的量纲为 4 的算符为
\begin{align}
  (A_1)^a_\mu(x) &= \frac{1}{2}\,\bar\psi(x)\gamma_5\sigma_{\mu\nu}
    \left[\overrightarrow{D}_\nu(x) - \overleftarrow{D}_\nu(x)\right]
    \tau^a\psi(x) , \nonumber\\
  (A_2)^a_\mu(x) &= \frac{1}{2}\,\partial_\mu\left[\bar\psi(x)\gamma_5\tau^a\psi(x)\right] , \nonumber\\
  (A_3)^a_\mu(x) &= \frac{m}{2}\,\bar\psi(x)\gamma_\mu\gamma_5\tau^a\psi(x) .
\end{align}
与作用量的改进一样，我们可以应用场方程，用它来证明 $A_1$ 可以写成 $A_2$ 和 $A_3$ 的线性组合。此外，$A_3$ 具有 (9.13) 中原始流的形式，因此它可以被吸收进轴矢流的总乘法重整化常数 $Z_A$ 之中（比较第 11 章）。于是，对于 $\order{a}$ 改进的轴矢流 $A_I$，我们得到如下表达式
\begin{equation}
  (A_I)^a_\mu(n) = A^a_\mu(n) + c_A\,a\,\widetilde{\partial}_\mu\,P^a(n) ,
\end{equation}
其中实系数 $c_A$ 仍有待确定。我们用 $\widetilde{\partial}_\mu$ 表示对称差分算符 $\widetilde{\partial}_\mu f(n) = (f(n + \hat\mu) - f(n - \hat\mu))/2a$。

对于赝标量密度 $P^a$，出现的唯一项是带额外质量因子的归一化，这同样被吸收进重整化之中。因此，对于 $\order{a}$ 改进，使用 $(P_I)^a = P^a$ 就够了。

\subsection{改进系数的确定}
\label{sec:9.1.4}

为了实现轴矢流 $n$ 点函数的 $\order{a}$ 改进，我们仍需要确定系数 $c_{sw}$ 和 $c_A$。如已指出的那样，在这里 QCD 由于它的非线性以及必要的重整化而偏离我们的玩具例子。

利用传统方法，可以在微扰论中计算这些系数 [5, 10, 11]：
\begin{equation}
  c_{sw} = 1 + 0.2659\,g^2 + O(g^4) , \qquad c_A = -0.00756\,g^2 + O(g^4) .
\end{equation}
然而，格点 QCD 的许多重要应用所处的区域中耦合 $g$ 并不小，微扰结果并不特别有用。另一种方法（即蝌蚪图改进 [12]）的入门性讨论见 [13]。

我们已经指出，对于 $c_{sw}$ 和 $c_A$ 的非微扰确定，可以利用 QCD 的手征对称性 [6]。在第 11 章中我们将证明，可以导出真空期望值之间的非平凡关系，即所谓的沃德恒等式。这里我们只陈述它们的存在而不加证明。我们在此使用的特定沃德恒等式在连续情形下为
\begin{equation}
  \ev{\partial_\mu A^a_\mu(x)\,O} = 2m\,\ev{P^a(x)\,O} .
\end{equation}
在这个方程中，$A^a_\mu$ 和 $P^a$ 是 (9.13) 中定义的相应格点插值算符的连续对应物。$O$ 是选取得使 $\ev{P^a(x)\,O}$ 不消失的插值算符，并限制其支集不包含 $A^a_\mu(x)$ 和 $P^a(x)$ 所在的点 $x$。方程 (9.17) 表达了这样一个事实：对于 $m = 0$ 的手征理论，同位旋矢量轴矢流是守恒的。因此它被称为部分守恒轴矢流关系（PCAC）。

在格点上，我们预期在重整化之后（见第 11 章），格点插值算符趋近连续关系 (9.17)，使得
\begin{equation}
  \ev{\widetilde{\partial}_\mu A^{(r)a}_\mu(n)\,O}
  = 2m^{(r)}\,\ev{P^{(r)a}(n)\,O} + O(a^z) .
\end{equation}
该方程中的上标 $(r)$ 表示使用的是重整化量。重要的观察是，对于适当选取了 $c_{sw}$ 和 $c_A$ 的改进理论，修正 $O(a^z)$ 是不同的，此时预期 $z = 2$；而对于未改进的理论，预期有 $\order{a}$ 修正，即 $z = 1$。因此，我们可以利用 PCAC 关系 (9.18) 中的修正来确定 $c_{sw}$ 和 $c_A$。

必要的重整化带来了一个小复杂性。使用改进的轴矢流 $(A_I)^a_\mu$（对于赝标量我们使用原始定义 $(P_I)^a = P^a$），我们把重整化插值算符写为
\begin{align}
  A^{(r)a}_\mu(n) &= Z_A(1 + b_A\,a\,m)\left[A^a_\mu(n) + c_A\,a\,\widetilde{\partial}_\mu P^a(n)\right] , \nonumber\\
  P^{(r)a}(n) &= Z_P(1 + b_P\,a\,m)\,P^a(n) .
\end{align}
这里 $Z_A$ 和 $Z_P$ 分别是轴矢流和赝标量的重整化常数（见第 11 章）。实系数 $b_A$ 和 $b_P$ 参数化它们的格点修正，包括诸如 $(A_3)^a_\mu$ 之类的修正项，这些项如我们已指出的那样被吸收进重整化之中。

基于 (9.18)，我们通过如下比值定义一个未重整化夸克质量 $m_{\mathrm{AWI}}$，称为 PCAC 质量或 AWI（轴矢沃德恒等式）质量
\begin{equation}
  m_{\mathrm{AWI}}(n)
  = \frac{\left\langle \left[\widetilde{\partial}_\mu A^a_\mu(n) + c_A\,a\,\Delta P^a(n)\right] O^a \right\rangle}
         {2\left\langle P^a(n)\,O^a \right\rangle} .
\end{equation}
在这个方程中，$O^a$ 是我们将在下面定义的插值算符，$\Delta$ 表示格点上的拉普拉斯算符，$\Delta f(n) = \sum_\mu (f(n + \hat\mu) - 2f(n) + f(n - \hat\mu))/a^2$。由格点 PCAC 关系 (9.18) 我们得到
\begin{equation}
  m_{\mathrm{AWI}} = \frac{Z_P(1 + b_P\,a\,m)}{Z_A(1 + b_A\,a\,m)}\,m^{(r)} + O(a^z) ,
\end{equation}
其中当 $c_{sw}$ 和 $c_A$ 具有正确的非微扰值时，得到 $z = 2$。为了确定这些值，在 (9.20) 中使用三个不同的 $O^a$ 对 $m_{\mathrm{AWI}}$ 求值三次。这是在固定的裸参数下进行的，使得三个值中的一个可用于消去比值 $(Z_P(1+b_P am))^{-1}(Z_A(1+b_A am))$，另外两个则确定 $c_{sw}$ 和 $c_A$。

剩下的问题是插入 (9.20) 的算符 $O^a$ 的选择。算符 $O^a$ 必须具有 $P^a$ 的量子数，并且应当选择得使它们能探测 PCAC 关系（从而探测截断效应），但同时保持 (9.21) 中的重整化因子不变。具有这些性质的算符 [6] 可以利用所谓的薛定谔泛函来构造。

对于薛定谔泛函 [14--18]，在时间方向上不使用周期（或反周期）边界条件，而是使用狄利克雷，即固定边界条件。空间方向保持周期边界条件。使用一个在时间方向上有 $N_T + 1$ 个格点的格子，我们需要指定 $n_4 = 0$ 和 $n_4 = N_T$ 处的场边界值：
\begin{align}
  P_+\psi(n,0) &= \rho(n) , & \bar\psi(n,0)P_- &= \bar\rho(n) , \nonumber\\
  P_-\psi(n,N_T) &= \rho'(n) , & \bar\psi(n,N_T)P_+ &= \bar\rho'(n) ,
\end{align}
其中 $P_\pm$ 是投影算符 $P_\pm = (1 \pm \gamma_4)/2$。对于规范场，我们需要指定 $n_4 = 0$ 和 $n_4 = N_T$ 处空间规范链的值：
\begin{align}
  U_j(n,0) &= \mathrm{diag}\left(e^{i\varphi_1/N}, e^{i\varphi_2/N}, e^{i\varphi_3/N}\right) , \nonumber\\
  U_j(n,N_T) &= \mathrm{diag}\left(e^{i\varphi'_1/N}, e^{i\varphi'_2/N}, e^{i\varphi'_3/N}\right) .
\end{align}
这里 $N$ 表示空间方向的格点数目，相位 $\varphi_i$、$\varphi'_i$ 满足条件 $\varphi_1 + \varphi_2 + \varphi_3 = 0$ 和 $\varphi'_1 + \varphi'_2 + \varphi'_3 = 0$。

在构成下述 (9.25) 中薛定谔泛函的路径积分中被积掉的动力学变量，正是我们格子内部的费米子和规范自由度。另一方面，边界值可以用作源项。利用这一特性，我们写下插入 (9.21) 的插值算符 $O^a$，
\begin{equation}
  O^a = -\frac{a^6}{2}\sum_{n,m}\gamma_5\tau^a
    \frac{\partial}{\partial\rho(n)}\frac{\partial}{\partial\bar\rho(m)} .
\end{equation}
于是，例如 $A^a_\mu$ 与 $O^a$ 的关联子在薛定谔泛函中的真空期望值定义为
\begin{equation}
  \ev{A^a_\mu(n)\,O^a}
  = \frac{1}{Z}\int D[\psi,\bar\psi,U]\,A^a_\mu(n)\,O^a\,e^{-S[\psi,\bar\psi,U]}
    \Big|_{\rho=\bar\rho=\rho'=\bar\rho'=0} ,
\end{equation}
即在完成 $O^a$ 定义中关于 $\rho$、$\rho'$、$\bar\rho$ 和 $\bar\rho'$ 的求导之后，把费米子的边界值置为 0。规范场的边界值保持在非平凡的值，现在可用于探测系统。

算符 $O^a$ 对应于放置在时间切片 $n_4 = 0$ 处、空间动量为零的夸克源和汇。因此真空期望值 (9.25) 可以解释为一条夸克线和一条反夸克线行进到我们格子的内部，在点 $n$ 处湮灭，即我们放置流 $A^a_\mu(n)$ 的点。

对我们的应用而言，薛定谔泛函的关键性质是，在 (9.21) 中重整化常数和重整化质量不依赖于边界条件 $\varphi_i$ 和 $\varphi'_i$。因此它们可以用来探测参数化离散化误差的指数 $z$。如果 $c_{sw}$ 和 $c_A$ 选择得当，指数为 $z = 2$，并且如 (9.20) 中定义的 PCAC 质量 $m_{\mathrm{AWI}}$ 的结果应该（基本上）独立于 $\varphi_i$、$\varphi'_i$ 以及插入流 $A^a_\mu$ 的坐标 $n_4$。关于从这一条件确定 $c_{sw}$ 和 $c_A$ 的策略的技术细节，我们请读者参阅原始文献 [19--21]。

两味动力学 Wilson 费米子情形下系数 $c_{sw}$ 的数值结果可以参数化为 [21]
\begin{equation}
  c_{sw} = \frac{1 - 0.454g^2 - 0.175g^4 + 0.012g^6 + 0.045g^8}{1 - 0.720g^2} ,
\end{equation}
在耦合 $g \in [0.0, 1.1]$ 的范围内这是可靠的近似。

\section{由重整化群变换得到的自由费米子格点作用量}
\label{sec:9.2}

在上一节中，我们通过分析格点作用量离散化中可能出现的不同项如何依赖于格距 $a$ 来实现改进。我们看到，以适当选取的系数对这些项作线性组合可以减小离散化误差。在接下来的三节中，我们讨论一种不同的方法，即所谓的重整化群（RG）变换，其中格点截断以上的自由度被积掉，这些紫外模的效应被纳入格点作用量的适当项中。

在本节中，我们直接从连续理论阻塞自由费米子（见 [22]）。这一热身练习通过一个可以解析求解的例子展示了 RG 变换的一些核心概念。这个例子对于格点上手征费米子问题也很有意义。我们将看到，由阻塞过程得到的格点作用量满足 Ginsparg--Wilson 方程，我们在第 7 章已指出它是格点手征对称性的关键。

\subsection{在超立方体上积掉场}
\label{sec:9.2.1}

出发点是自由无质量费米子的连续作用量：
\begin{equation}
  S_F[\varphi,\bar\varphi] = \int d^4x\,\bar\varphi(x)\,\gamma_\mu\partial_\mu\,\varphi(x)
  = \int d^4p\,\bar\varphi(-p)\,i\gamma_\mu p_\mu\,\varphi(p) .
\end{equation}
为便于后文使用，我们在第二步中已经用对称约定
\begin{equation}
  \bar f(p) = \int_{-\infty}^{\infty}\frac{d^4x}{(2\pi)^2}\,f(x)\,e^{-ip\cdot x} ,
  \qquad
  f(x) = \int_{-\infty}^{\infty}\frac{d^4p}{(2\pi)^2}\,\bar f(p)\,e^{ip\cdot x}
\end{equation}
把作用量变换到动量空间。我们在此强调，连续理论中的傅里叶变换用 $\bar f$ 表示，而格点上的傅里叶变换用 $\hat f$ 表示，如附录 A.3 中所定义。

连续场 $\varphi(x)$、$\bar\varphi(x)$ 只携带狄拉克指标，因为没有耦合规范场。为了从连续场 $\varphi(x)$、$\bar\varphi(x)$ 得到阻塞的格点场 $\varphi^B_n$、$\bar\varphi^B_n$，我们把一个无穷欧几里得格子 $\mathbb{Z}^4$ 嵌入到我们的欧几里得时空 $\mathbb{R}^4$ 中。格常数 $a$ 暂时取为 $a = 1$，但稍后在需要时会显式写出。我们通过把连续场在以格子格点 $n$（即 $n \in \mathbb{Z}^4$）为中心、边长为 1 的超立方体上积分来构造阻塞场：
\begin{equation}
  \varphi^B_n = \int_{-\infty}^{\infty}\frac{d^4p}{(2\pi)^2}
    \bar\varphi(p)\,\pi(p)\,e^{ip\cdot n} ,
  \qquad
  \pi(p) = \prod_{\mu=1}^4\int_{n_\mu-\frac{1}{2}}^{n_\mu+\frac{1}{2}}dx_\mu\,e^{ip\cdot x} .
\end{equation}
在这个方程中，我们把连续场 $\varphi(x)$ 用 (9.28) 中定义的傅里叶变换 $\bar\varphi(p)$ 表示，并把剩余的时空积分缩写为
\begin{equation}
  \pi(p) = \prod_{\mu=1}^4\int_{-\frac{1}{2}}^{\frac{1}{2}}dx_\mu\,e^{ip\cdot x}
  = \prod_{\mu=1}^4\frac{2\sin(p_\mu/2)}{p_\mu} .
\end{equation}
函数 $\pi(p)$ 描述动量空间中的阻塞过程。共轭场 $\bar\varphi$ 的阻塞以完全相同的方式进行。由于阻塞过程在狄拉克空间中是平凡的，阻塞场只是继承其连续对应物的狄拉克指标。

方程 (9.29) 把阻塞场 $\varphi^B_n$ 表示为用连续场的傅里叶变换 $\bar\varphi(p)$ 表示的傅里叶积分。为便于后文使用，我们改写 (9.29) 得到另一个表达式，其中使用阻塞场的格点傅里叶变换 $\hat\varphi^B(p)$：
\begin{equation}
  \varphi^B_n = \int_{-\pi}^{\pi}\frac{d^4p}{(2\pi)^4}\,\hat\varphi^B(p)\,e^{ip\cdot n} ,
  \qquad
  \hat\varphi^B(p) = \sum_{k\in\mathbb{Z}^4}\bar\varphi(p+2\pi k)\,\pi(p+2\pi k) .
\end{equation}
注意，现在动量积分只跑遍布里渊区，即 $p_\mu \in (-\pi, \pi]$，这是无穷格子的动量空间，从附录 A.3 中的公式可以看出，只需把格点数目推向无穷即可。

\subsection{阻塞的格点狄拉克算符}
\label{sec:9.2.2}

在定义了阻塞场 $\varphi^B_n$、$\bar\varphi^B_n$ 并把它们的表达式化为傅里叶积分之后，我们现在可以形式地定义由阻塞过程得到的格点理论：
\begin{multline}
  \exp\!\left[-\sum_{n,m}\bar\psi_n\,D_{nm}\,\psi_m\right]
  = \int D[\varphi,\bar\varphi]\,
    \exp\!\left[-\sum_n\left(\bar\psi_n-\bar\varphi^B_n\right)
      \left(\psi_n-\varphi^B_n\right)\right. \\
  \qquad\qquad \left.- S_F[\varphi,\bar\varphi]\right] .
\end{multline}
在这个方程中，$\psi_m$、$\bar\psi_n$ 表示目标格点理论的场，即生活在无穷格子 $\mathbb{Z}^4$ 上的费米子。因此指标 $n$ 和 $m$ 跑遍整个 $\mathbb{Z}^4$。格点上的狄拉克算符记作 $D_{nm}$。左端于是就是格点理论的玻尔兹曼因子 $\exp(-S_{\mathrm{latt}})$。格点场 $\psi_m$、$\bar\psi_n$ 和格点狄拉克算符 $D_{nm}$ 都携带我们略去的旋量指标，即我们在狄拉克空间中使用矢量/矩阵记号。

(9.32) 的右端定义了如何从连续对应物得到格点理论，或者更准确地说得到格点狄拉克算符 $D_{nm}$：连续场 $\varphi$、$\bar\varphi$ 在路径积分中被积掉（我们稍后将评论连续测度 $D[\varphi,\bar\varphi]$）。在玻尔兹曼因子中，像通常一样有来自 (9.27) 的连续作用量 $S_F[\varphi,\bar\varphi]$，还有一个把阻塞连续场 $\varphi^B$、$\bar\varphi^B$ 与格点上目标理论的场 $\psi$、$\bar\psi$ 耦合起来的附加二次型部分。

由于右端的路径积分是高斯的，可以解析求解：方便的做法 [22] 是用对辅助格拉斯曼变量 $\eta$、$\bar\eta$ 的额外高斯积分来改写阻塞项：
\begin{multline}
  \exp\!\left[-\sum_n\left(\bar\psi_n-\bar\varphi^B_n\right)
    \left(\psi_n-\varphi^B_n\right)\right] \\
  = \int D[\eta,\bar\eta]\;
    \exp\!\left[\sum_n\left(\bar\eta_n\eta_n
      + \left(\bar\psi_n-\bar\varphi^B_n\right)\eta_n
      + \bar\eta_n\left(\psi_n-\varphi^B_n\right)\right)\right] .
\end{multline}
辅助场生活在格点上，携带一个狄拉克指标，测度 $D[\eta,\bar\eta]$ 是所有格点处测度的乘积。(9.33) 中混合辅助场与阻塞场 $\varphi^B$ 的项写为
\begin{align}
  \sum_n\bar\eta_n\,\varphi^B_n
  &= \int_{-\infty}^{\infty}\frac{d^4p}{(2\pi)^2}\,\bar J(-p)\,\bar\varphi(p) , \qquad
    \bar J(p) = \sum_n\bar\eta_n\,\pi(-p)\,e^{-ip\cdot n} , \nonumber\\
  \sum_n\bar\varphi^B_n\,\eta_n
  &= \int_{-\infty}^{\infty}\frac{d^4p}{(2\pi)^2}\,\bar\varphi(-p)\,J(p) , \qquad
    J(p) = \sum_n\eta_n\,\pi(-p)\,e^{-ip\cdot n} ,
\end{align}
其中我们使用了 (9.29) 通过连续傅里叶变换来表示阻塞场。

把 (9.33) 和 (9.34) 代入主公式 (9.32)，我们得到
\begin{equation}
  \exp\!\left[-\sum_{n,m}\bar\psi_n\,D_{nm}\,\psi_m\right]
  = \int D[\eta,\bar\eta]\,
    \exp\!\left[\sum_n\left(\bar\eta_n\eta_n + \bar\psi_n\eta_n + \bar\eta_n\psi_n\right)\right]
    I[\bar J,J] ,
\end{equation}
其中我们把对连续场剩余的路径积分记作 $I[\bar J,J]$，它由下式给出
\begin{equation}
  \begin{split}
    I[\bar J,J]
    &= \int D[\varphi,\bar\varphi]\,
      \exp\!\left[-\int_{-\infty}^{\infty}\frac{d^4p}{(2\pi)^2}\,
        \left(\bar\varphi(-p)\,i\gamma_\mu p_\mu\,\varphi(p)\right.\right. \\
    &\qquad\qquad\qquad
        \left.\left.+ \bar J(-p)\,\varphi(p) + \bar\varphi(p)\,J(p)\right)\right] \\
    &= C\,\exp\!\left[\int_{-\infty}^{\infty}\frac{d^4p}{(2\pi)^4}\,
        \bar J(-p)\,\frac{-i\gamma_\mu p_\mu}{p^2}\,J(p)\right] .
  \end{split}
\end{equation}
在第二步中，如第 5 章所讨论的，我们通过在指数中配平方并平移路径积分中的积分变量 $\varphi'(p) = \varphi(p) + i\gamma_\mu p_\mu\,(2\pi)^{-2}p^{-2}J(p)$（对 $\bar\varphi$ 类似）来求解路径积分。在积分变量平移后与 $\bar J$、$J$ 无关的剩余项记作常数 $C$。我们在此强调，我们并没有显式定义连续测度 $D[\varphi,\bar\varphi]$，而只是假设它在变量平移下保持不变。事实上，要适当定义这一测度，需要诸如有限格子之类的正规化子。如果去掉这一正规化子，常数 $C$ 会发散，但这对于算符真空期望值的计算无关紧要，因为在期望值用配分函数 $Z$ 归一化时这样的常数会消去。因此我们可以忽略常数 $C$，路径积分 (9.36) 就由配平方得到的 $J$ 的二次型项给出。

最后一步是把结果 (9.36) 代入对 $\eta$ 和 $\bar\eta$ 的剩余积分 (9.35)：
\begin{align}
  \exp\!\left[-\sum_{n,m}\bar\psi_n\,D_{nm}\,\psi_m\right]
  &= C\int D[\eta,\bar\eta]\;
     \exp\!\left[\sum_{n,m}\bar\eta_n\,(1+R)_{nm}\,\eta_m\right. \nonumber\\
  &\qquad \left.+ \sum_n\left(\bar\psi_n\eta_n + \bar\eta_n\psi_n\right)\right] \nonumber\\
  &= CC'\,\exp\!\left[-\sum_{n,m}\bar\psi_n\,(1+R)^{-1}_{nm}\,\psi_m\right] .
\end{align}
在最后一步中，我们再次通过在指数中配平方来求解高斯路径积分，与 $\psi$、$\bar\psi$ 无关的剩余项产生另一个无关紧要的常数 $C'$。矩阵 $R_{nm}$ 是通过把 (9.34) 中 $J(p)$ 和 $\bar J(p)$ 的表达式代入 (9.36) 右端的指数得到的。矩阵 $R$ 的显式为
\begin{equation}
  R_{nm} = \int_{-\infty}^{\infty}\frac{d^4p}{(2\pi)^2}\,
    \pi(-p)\,e^{-ip\cdot m}\,\frac{-i\gamma_\mu p_\mu}{p^2}\,\pi(p)\,e^{ip\cdot n} ,
\end{equation}
通过比较 (9.37) 中第一项和最后一项，我们得到格点狄拉克算符的最终结果（在略去常数 $C$ 和 $C'$ 之后）
\begin{equation}
  D_{nm} = (1+R)^{-1}_{nm} .
\end{equation}
格点狄拉克算符 $D$ 用其傅里叶变换来表示最为方便：
\begin{equation}
  \hat D(p|q) = \frac{1}{(2\pi)^4}\sum_{n,m}e^{-ip\cdot n}\,D_{nm}\,e^{iq\cdot m} .
\end{equation}
对于所需的 $R$ 的傅里叶变换，可以求得
\begin{align}
  \hat R(p|q)
  &= \delta(p-q)\sum_{k\in\mathbb{Z}^4}
     \frac{\pi(q+2\pi k)^2}{(q+2\pi k)^2}\,(-i)\gamma_\mu\,(q+2\pi k)_\mu \nonumber\\
  &= -i\,\delta(p-q)\,\gamma_\mu\,v_\mu(q) , \nonumber\\
  \text{with}\quad v_\mu(q)
  &= \sum_{k\in\mathbb{Z}^4}
     \frac{\pi(q+2\pi k)^2}{(q+2\pi k)^2}\,(q+2\pi k)_\mu .
\end{align}
把它代入 (9.39) 得到 [22]
\begin{equation}
  \hat D(p|q) = \delta(p-q)\,\hat D(q) , \qquad
  \hat D(q) = \left(a\,\mathbb{1} - i\gamma_\mu v_\mu(q)\right)^{-1}
    = \frac{a\,\mathbb{1} + i\gamma_\mu v_\mu(q)}{a^2 + v(q)^2} .
\end{equation}
这里我们分离出保证动量守恒的狄拉克 $\delta$，动量守恒是平移不变性的结果。我们还重新引入了下述分析所需的格距 $a$。

\subsection{阻塞作用量的性质}
\label{sec:9.2.3}

动量空间中的格点传播子 $\hat D(q)$ 是我们为得到有效格点理论而进行的阻塞过程的结果。它有一系列值得注意的性质，我们现在就来讨论。

首先我们确认 $\hat D(q)$ 具有正确的朴素连续极限，并且没有加倍体。换言之，我们需要证明
\begin{equation}
  \hat D(q) = \begin{cases}
    i\gamma_\mu q_\mu + O(a) & \text{when all } q_\mu \approx 0 ,\\
    O(1/a) & \text{when at least one } q_\mu \approx \pi/a .
  \end{cases}
\end{equation}
在审视格点狄拉克算符的最终结果 (9.42) 时，很明显性质 (9.43) 必定来自 (9.41) 中给出的函数 $v_\mu(q)$ 的行为。事实上，可以证明（例如通过数值分析 $v_\mu(q)$ 最快）它们的行为如下
\begin{equation}
  v_\mu(q) \approx \begin{cases}
    q_\mu/q^2 & \text{when all } q_\mu \approx 0 ,\\
    0 & \text{when at least one } q_\mu \approx \pi/a .
  \end{cases}
\end{equation}
这一行为恰好蕴含 (9.43) 中陈述的格点狄拉克算符所需性质。

在确认了阻塞狄拉克算符 $\hat D$ 给出一个合适的格点理论之后，让我们进一步研究它的性质。当变换回实空间时，发现阻塞狄拉克算符不是超局域的，这与第 7.4.2 节的一般讨论一致。然而，可以证明 [22] 阻塞狄拉克算符按指数衰减，这对于局域量子场论是必需的。

另一个重要方面是离散化误差。我们在 (9.43) 中看到，在动量 $q$ 趋于零的极限下，我们得到正确的行为 $i\gamma_\mu q_\mu$。然而，随着动量增大，与这一理想行为的偏差是什么呢？评估这个问题的物理上有趣的方式是分析由 $\hat D$ 描述的自由费米子的色散关系。这可以通过计算 $\hat D(q)^{-1}$ 在复 $q_4$ 平面中、对于空间动量 $q$ 的不同取值的极点来实现。把由极点位置 $q_4$ 的模给出的能量 $E(q)$ 对空间动量作图，可以发现对于阻塞狄拉克算符，这些值正好落在线性连续色散关系 $E(q) = |q|$ 之上，基本上一直到截断 $\pi/a$ 都是如此。这表明离散化误差很小。

最后我们评论阻塞格点狄拉克算符 $\hat D(q)$ 的手征性质。证明它满足如下形式的 Ginsparg--Wilson 方程是一个简单的练习
\begin{equation}
  \gamma_5\hat D(q) + \hat D(q)\gamma_5 = 2a\,\hat D(q)\gamma_5\hat D(q) .
\end{equation}
我们指出，这里遇到的 Ginsparg--Wilson 方程与原始版本 (7.29) 的归一化略有不同，即在右端多了一个因子 2。然而这个因子完全无关紧要，因为我们在第 7 章给出的格点手征对称性的整个纲领也可以用修改后的归一化来实现。因此我们发现阻塞理论也是手征的。

为总结本节，我们已经证明，从连续理论阻塞自由费米子会给出
\begin{itemize}
  \item 没有加倍体且具有局域作用量的合适的格点费米子。
  \item 几乎直到截断都具有完美色散关系的费米子。
  \item 满足 Ginsparg--Wilson 方程的狄拉克算符。
\end{itemize}
有了这三个性质，很明显阻塞过程是构造量子场论格点离散化的强有力方法。在第 9.4 节中我们将证明，涌现出的手征格点对称性并非巧合，而是连续对称性到由阻塞得到的相应格点理论的一个相当一般映射的结果。然而，在进入这一分析之前，我们需要处理如何把这里为自由费米子概述的策略推广到胶子与夸克场相互作用的完整 QCD 情形。

\section{QCD 的实空间重整化群}
\label{sec:9.3}

在上一节中，我们通过对自由费米子的显式计算证明了，积掉截断 $\sim 1/a$ 以上自由度的 RG 方法是构造量子场论格点离散化的强有力工具。所得的格点作用量具有优美的性质，特别是离散化效应几乎一直到截断都不存在，而且格点狄拉克算符满足 Ginsparg--Wilson 方程，这是连续手征对称性在格点上的体现。

对于完整 QCD 的情形，即规范场与费米子耦合时，情况要复杂得多。一个严重的障碍阻碍了像上一节考虑的自由费米子那样直接从连续理论进行阻塞：(9.32) 右端的连续路径积分在出现非高斯项时无法解析求解。这样的项有规范场与费米子的相互作用项，以及规范作用量的三次和四次项。

概念上的出路是在一个非常细的格子上使用 QCD 的格点版本，即一个接近连续极限、格距 $a$ 很小且离散化效应微小的格子。这个细格子取代连续理论作为阻塞变换的出发点。我们不是积分连续理论的超立方体，而是在细格子的一个格点块上对格点场求平均，以构造更粗格子上的场。粗格子上的作用量通过一个等价于 (9.32) 的方程来识别。迭代这样的实空间重整化群变换，我们可以得到目标理论在所需粗细的格子上的作用量。直到截断，这个粗作用量都与我们出发的细格子上的理论具有相同的离散化误差。在 RG 变换中被平均掉的短距离紫外模的物理，被编码在粗格子所用作用量的各项之中。我们指出，实空间重整化群最初是作为分析接近临界点的统计力学系统的方法而发展起来的 [23, 24]。关于它在格点 QCD 中的应用、超出我们这里简短描述的入门介绍，我们推荐 [25]。

\begin{figure}[htbp]
  \centering
  \includegraphics[width=0.85\textwidth]{fig92.png}
  \caption*{Fig. 9.2. \emph{重叠离散块自旋变换的一个例子。我们在格子的一个平面中展示了如何对细格子上的费米子场（实心点）求平均，以构造生活在粗格子上的阻塞费米子场（围绕只有偶指标的格点的空心圆）。沿箭头方向包含规范输运子，以在粗格子上得到规范协变的旋量}.}
  \label{fig:92}
\end{figure}

\subsection{阻塞完整的 QCD}
\label{sec:9.3.1}

让我们从描述如何从格常数为 $a$ 的细格子 $\Lambda$ 上的场构造生活在格距为 $a' = 2a$ 的粗格子 $\Lambda'$ 上的新场开始。由于我们现在考虑 QCD，必须同时阻塞费米子和规范场。在讨论实际构造之前，我们设定记号约定：细格子 $\Lambda$ 的时空指标记作 $n$，粗格子 $\Lambda'$ 的时空指标记作 $n'$。细格子上的费米子和规范场为 $\psi_n$、$\bar\psi_n$ 和 $U_\mu(n)$，它们生活在粗格子上的阻塞对应物为 $\psi^B_{n'}$、$\bar\psi^B_{n'}$ 和 $U^B_\mu(n')$，真正的粗场记作 $\psi'_{n'}$、$\bar\psi'_{n'}$ 和 $U'_\mu(n')$。

我们从讨论费米子的阻塞方案开始。许多不同的阻塞方案是可能的。图 9.2 展示了一个相当简单的方案，图中我们展示了如何在格子的一个平面中阻塞场（其他平面等价地阻塞）。费米子场 $\psi_n$、$\bar\psi_n$ 生活在细格子的格点上（图中实心点），需要适当地求平均，以得到生活在粗格子上的阻塞费米子 $\psi^B_{n'}$、$\bar\psi^B_{n'}$，我们选取的粗格子由细格子中只有偶指标的格点构成（图中以圆圈标出）。粗格子格点 $n'$ 处的阻塞场 $\psi^B_{n'}$、$\bar\psi^B_{n'}$ 通过对所有与 $n'$ 由一条或多条定向路径相连的格点 $n$ 上的场 $\psi_n$、$\bar\psi_n$ 求平均得到。显然，细格子的格点与粗格子的多个格点相连。这种类型的阻塞方案被称为重叠块自旋变换。

在组合来自细格子不同格点 $n$ 的费米子场 $\psi_n$、$\bar\psi_n$ 时，必须保持规范协变性完好。因此我们需要通过规范输运子 $U(n',n)$ 把格点 $n$ 处的场连接到放置阻塞场的格点 $n'$。这样的输运子通过沿图中定向路径相乘细格子的链变量 $U_\mu(m)$ 得到。如果一条链沿负方向穿过，我们应用通常的约定 $U_{-\mu}(m) \equiv U_\mu(m-\hat\mu)^\dagger$。为了得到尊重格子对称性的阻塞方案，几个格点 $n$ 通过一组由离散转动、反射等对称性联系起来的路径连接到 $n'$。我们可以把费米子的阻塞总结为
\begin{equation}
  \psi^B_{n'} = \sum_{n\in\mathcal{N}(n')}
    \left(\sum_{p\in\mathcal{P}(n',n)}s_p\prod_{(m,\mu)\in p}U_\mu(m)\right)\psi_n
  \equiv \sum_n\omega[U]_{n'n}\,\psi_n .
\end{equation}
这里 $\mathcal{N}(n')$ 是细格子中围绕粗格子上目标格点 $n'$ 的格点 $n$ 的一个邻域。用 $\mathcal{P}(n',n)$ 表示我们用来连接细格子格点 $n$ 与 $n'$ 的路径 $p$ 的集合，例如图 9.2 中描绘的选择。乘积跑遍路径 $p$ 中所有链 $(m,\mu)$。每条路径带有一个系数 $s_p$。这些系数必须反映对称性（例如由转动联系起来的路径必须有相同的系数），并且作为归一化条件，我们要求连接到目标格点 $n'$ 的所有路径的系数之和为 1。在 (9.46) 的最后一步中，我们通过引入一个依赖于细格子规范场 $U$、收集连接 $n'$ 与 $n$ 的所有项的矩阵 $\omega[U]_{n'n}$ 来改写阻塞方案。共轭场 $\bar\psi(n)$ 用 $\omega[U]$ 的厄米共轭等价地阻塞。我们指出，与上一节一样，阻塞过程在狄拉克空间中是平凡的。

关于规范链 $U_\mu(n)$ 的阻塞，一种简单的做法是把连接粗格子两个格点的两条细链变量相乘，从而得到阻塞变量。然而，这里通常也使用更一般的规范协变平均，一个简单的版本由下式给出
\begin{align}
  U^B_\mu(n') &= (1-\alpha)\,U_\mu(n')\,U_\mu(n'+\hat\mu) \nonumber\\
    &\quad + \frac{\alpha}{6}\sum_{\nu\neq\mu}
      \left[U_\nu(n')\,U_\mu(n'+\hat\nu)\,U_\mu(n'+\hat\nu+\hat\mu)\,
       U_\nu(n'+2\hat\mu)^\dagger \right. \nonumber\\
    &\qquad \left. + U_\nu(n'-\hat\nu)^\dagger\,U_\mu(n'-\hat\nu)\,
       U_\mu(n'-\hat\nu+\hat\mu)\,U_\nu(n'-\hat\nu+2\hat\mu)\right] .
\end{align}
这里 $\hat\mu$ 和 $\hat\nu$ 是细格子上指向 $\mu$ 和 $\nu$ 方向的单位矢量，$\alpha$ 是阻塞过程的实参数。本质上，阻塞对细格子上连接粗格子相邻格点的链路径作线性组合。沿这些路径考虑细链变量的有序乘积。找到最优的阻塞方案需要一定的实验摸索，比 (9.47) 中的方案更精细的方案在 [26--29] 中讨论。

我们指出，在上述定义中，阻塞场 $U^B_\mu(n')$ 不是规范群的元素，因为两个 SU(3) 矩阵之和不在 SU(3) 中。下面我们将定义阻塞场 $U^B_\mu(n')$ 如何与粗格子上的链变量 $U'_\mu(n')$ 相联系，后者当然将是规范群的元素。

在定义了如何得到阻塞场 $\psi^B$、$\bar\psi^B$ 和 $U^B$ 之后，我们现在可以表述一个推广 (9.32)、定义粗格子上理论的阻塞关系：
\begin{multline}
  e^{-S'_F[\psi',\bar\psi',U'] - \beta' S'_G[U']}
  = \int D[\psi,\bar\psi]\,D[U]\,e^{-S_F[\psi,\bar\psi,U] - \beta S_G[U]} \\
  \times e^{-T_F[\psi',\bar\psi',\psi^B,\bar\psi^B,U] - \beta\,T_G[U',U^B]} .
\end{multline}
在这个方程中，$S_F[\psi,\bar\psi,U]$ 和 $S'_F[\psi',\bar\psi',U']$ 是细格子和粗格子的费米子作用量，$\beta S_G[U]$ 和 $\beta' S'_G[U']$ 是相应的规范作用量，其中 $\beta$、$\beta'$ 表示两个格子的逆规范耦合。右端现在是一个对细格子上自由度的良定义的格点路径积分。阻塞函数 $T_F$ 和 $T_G$ 把粗格子上目标理论的费米子场和规范场 $\psi'$、$\bar\psi'$、$U'$ 连接到阻塞量 $\psi^B$、$\bar\psi^B$、$U^B$。对于费米子我们使用
\begin{align}
  T_F[\psi',\bar\psi',\psi^B,\bar\psi^B,U]
  &= \kappa_F\sum_{n'}\left(\bar\psi'_{n'} - b\,\bar\psi^B_{n'}\right)
     \left(\psi'_{n'} - b\,\psi^B_{n'}\right) \nonumber\\
  &= \kappa_F\sum_{n'}\left(\bar\psi'_{n'} - b\sum_n\bar\psi_n\,\omega[U]^\dagger_{nn'}\right) \nonumber\\
  &\qquad \times\left(\psi'_{n'} - b\sum_n\omega[U]_{n'n}\,\psi_n\right) ,
\end{align}
其中 $b$ 是实值归一化常数，$\kappa_F$ 是可用于调节粗格子上作用量性质（尤其是其局域性）的实参数。在 (9.49) 的第二行中，我们用 (9.46) 中引入的变换矩阵 $\omega[U]$ 表示阻塞场。对于规范场我们定义
\begin{align}
  T_G[U',U^B] &= \frac{\kappa_G}{3}\sum_{n',\nu}
    \left[\mathrm{Re}\,\tr\left(U'_\nu(n')\,U^B_\nu(n')^\dagger\right)
      - \mathcal{N}\left(U^B_\nu(n')\right)\right] , \nonumber\\
  \exp\left[\mathcal{N}\left(U^B_\nu(n')\right)\right]
  &\equiv \int dW\,\exp\left[\beta\kappa_G\,\mathrm{Re}\,\tr\left(W\,U^B_\nu(n')^\dagger\right)/3\right] .
\end{align}
第二个方程通过对 $W \in \mathrm{SU}(3)$ 的积分定义归一化 $\mathcal{N}(U^B_\nu(n'))$，$\kappa_G$ 同样是变换的实参数。

\subsection{耦合的 RG 流}
\label{sec:9.3.2}

在上一节中，我们定义了费米子和规范场的阻塞方案，从而给了联系粗格子和细格子上理论作用量的重整化群方程 (9.48) 一个良定义的含义。对于任意的粗位形 $\psi'$、$\bar\psi'$、$U'$，(9.48) 的右端在左端的指数中定义了粗格子上的作用量。

我们已经指出，RG 变换的关键思想是把细格子上模的物理包含在粗格子所用作用量的适当项中。这样的项例如可以像本章第一节引入的三叶草项 (9.10)。粗格子上作用量的所有这些项都带有某些系数 $c'_1$、$c'_2$、$\ldots$，对于三叶草项的例子，这些系数对应于 Symanzik 改进纲领的 Sheikholeslami--Wohlert 系数。我们假定可以把 RG 变换 (9.48) 看作参数化格点 QCD 作用量的耦合空间中的映射，
\begin{equation}
  \beta,\,c_1,\,c_2,\,\ldots\;\xrightarrow{\;\mathrm{RG}\;}
    \beta',\,c'_1,\,c'_2,\,\ldots ,
\end{equation}
其中 $\beta$、$c_i$ 是细格子上的耦合（对于可能是我们出发点的普通 Wilson 作用量，其中大部分为零）。为了标记耦合，我们不区分作用量的费米子和规范场部分，只单独标出逆规范耦合 $\beta$、$\beta'$。

从耦合为 $\beta$、$c_1$、$c_2$、$\ldots$ 的理论出发，我们阻塞场，粗格子的耦合 $\beta'$、$c'_1$、$c'_2$、$\ldots$ 取使得细格子被积掉的场的影响得以考虑的值。与此同时，格距以标度因子 2 增大，
\begin{equation}
  a \xrightarrow{\;\mathrm{RG}\;} a' = 2a .
\end{equation}
我们指出，标度因子 2 是图 9.2 和 (9.47) 所定义阻塞变换所特有的，其他阻塞方案可以有不同的标度因子。

我们现在问自己，映射 (9.51) 是否具有不动点（FP），即一组在 RG 变换 (9.51) 下映射到自身的耦合。方程 (9.52) 提示可能存在两种类型的不动点，分别对应于 $a = a' = \infty$ 或 $a = a' = 0$。显然只有第二种可能性（通常称为临界不动点）对我们有意义，因为它对应于连续极限 $a \to 0$。正如我们在第 3.5.4 节中所讨论的，格点 QCD 的连续极限（其中 $a \to 0$）在 $\beta \to \infty$ 的极限下得到。因此我们预期在 $\beta = \infty$ 和参数值 $c_1$、$c_2$、$\ldots$ 处找到合适的临界不动点。在我们的 RG 变换 (9.51) 下，这些值映射到自身。

在实际应用中，我们不能从 $\beta = \infty$ 开始进行 RG 变换。然而，对重复 RG 变换下耦合流的详细分析（例如见 [25]）表明，从一个大但有限的 $\beta$ 值出发就足够了，重复 RG 变换的少数几步将把耦合 $c_1$、$c_2$、$\ldots$ 驱动到不动点值 $c_1$、$c_2$、$\ldots$。

\subsection{RG 方程的鞍点分析}
\label{sec:9.3.3}

到目前为止，我们已经为完整 QCD 指定了阻塞方案，定义了 RG 变换 (9.48)，并识别了一个对应于我们所关心的连续物理的可能不动点。剩下要做的是找到求解 RG 变换 (9.48) 的方法，以便我们能在格子上构造 RG 作用量。

为此，我们可以利用想要找到的临界不动点对应于 $\beta = \beta' = \infty$ 这一事实。在这个极限下，(9.48) 右端的路径积分由其鞍点主导，即在细格子上使指数取极小值的规范场 $U$ 位形（不一定唯一）。(9.48) 两端指数中与逆耦合 $\beta$、$\beta'$ 相乘的项，在把逆耦合推向无穷时必须匹配。我们读出作用量规范场部分的如下方程：
\begin{equation}
  S'_G[U'] = \min_U\left\{S_G[U] + T^\infty_G[U',U^B]\right\} .
\end{equation}
附加在阻塞函数 $T_G$ 上的上标 $\infty$ 表示它也必须在 $\beta = \infty$ 处求值。由于我们想识别不动点处的所谓不动点（FP）作用量 $S^{\mathrm{FP}}_G$，它由映射到自身的系数 $c_i$ 定义，我们可以令 $S'_G = S_G = S^{\mathrm{FP}}_G$，得到识别不动点作用量规范部分 $S^{\mathrm{FP}}_G$ 的方程，
\begin{align}
  S^{\mathrm{FP}}_G[U']
  &= \min_U\left\{S^{\mathrm{FP}}_G[U] + T^\infty_G[U',U^B]\right\} \nonumber\\
  &= \min_U\Biggl\{S^{\mathrm{FP}}_G[U] + \frac{\kappa_G}{3}\sum_{n',\nu}
     \left[\mathrm{Re}\,\tr\left(U'_\nu(n')\,U^B_\nu(n')^\dagger\right)\right. \nonumber\\
  &\qquad \left.- \mathcal{N}^\infty\left(U^B_\nu(n')\right)\right]\Biggr\} ,
\end{align}
其中在第二步中我们插入了 (9.50) 中给出的阻塞函数 $T_G$ 的显式表达式。所需的极限 $\beta \to \infty$ 很容易实现，因为 $\beta$ 只通过归一化常数 $\mathcal{N}^\infty$ 进入，而后者又可以通过其鞍点求值给出
\begin{equation}
  \mathcal{N}^\infty\left(U^B_\nu(n')\right)
  = \frac{\kappa_G}{3}\,\max_{W\in\mathrm{SU}(3)}
    \mathrm{Re}\,\tr\left(W\,U^B_\nu(n')^\dagger\right) .
\end{equation}
在 (9.48) 的鞍点分析中匹配带因子 $\beta$、$\beta'$ 的指数，得到确定不动点作用量规范部分的方程 (9.54)。要求费米子贡献的其余因子也匹配，给出
\begin{equation}
  e^{-S'_F[\psi',\bar\psi',U']} = \int D[\psi,\bar\psi]\,
    e^{-S_F[\psi,\bar\psi,U^{\min}] - T_F[\psi',\bar\psi',\psi^B,\bar\psi^B,U^{\min}]} ,
\end{equation}
其中右端必须对使规范场 RG 方程 (9.54) 右端取极小值的位形 $U^{\min}$ 求值。

(9.56) 右端的积分可以求解，因为它是高斯的。借助粗格子和细格子的狄拉克算符 $D'$ 和 $D$，我们把粗格子和细格子的费米子作用量写为
\begin{align}
  S'_F[\psi',\bar\psi',U'] &= \sum_{n',m'}\bar\psi'_{n'}\,D'[U']_{n'm'}\,\psi'_{m'} , \nonumber\\
  S_F[\psi,\bar\psi,U^{\min}] &= \sum_{n,m}\bar\psi_n\,D[U^{\min}]_{nm}\,\psi_m .
\end{align}
利用阻塞函数 $T_F(\psi',\bar\psi',\psi^B,\bar\psi^B, U^{\min})$ 也是 $\psi$ 和 $\bar\psi$ 的双线性形式（见 (9.49)）这一事实，我们用第 5.1.4 节的技术求解费米子积分。(9.56) 右端积分的结果是粗场 $\psi'$、$\bar\psi'$ 的双线性形式的指数。这个双线性形式的核必须与 (9.56) 左端双线性形式指数中的核匹配。我们得到
\begin{equation}
  \begin{split}
    D'[U']_{n'm'}
    &= \kappa_F\,\delta_{n'm'}
       - \kappa_F^2 b^2\left[\omega[U^{\min}]
         \left(D[U^{\min}] + \kappa_F b^2\,\omega[U^{\min}]^\dagger
         \omega[U^{\min}]\right)^{-1}\right. \\
    &\qquad\qquad \left.\cdot \omega[U^{\min}]^\dagger\right]_{n'm'} ,
  \end{split}
\end{equation}
这就是把粗格子上的狄拉克算符 $D'$ 与细格子的狄拉克算符 $D$ 联系起来的 RG 方程。

方程 (9.54) 和 (9.58) 是联系细格子上作用量与粗格子上作用量的两个 RG 方程。有趣的是，不动点规范作用量的 RG 方程独立于费米子问题。然而，费米子的 RG 方程 (9.58) 需要细格子上的规范位形 $U^{\min}$ 作为输入，它使 (9.54) 的右端取极小值。在下一节中，我们讨论如何近似 (9.54) 和 (9.58) 的解。

\subsection{求解 RG 方程}
\label{sec:9.3.4}

这两个 RG 方程 (9.54) 和 (9.58) 是非线性方程，应当对任意的粗背景规范场 $U'$ 成立。第一步必须确定与给定的 $U'$ 对应并使 (9.54) 右端取极小值的细位形 $U^{\min}$。然后这个位形被用于费米子的 RG 方程 (9.58)。

为了确定不动点规范作用量 $S^{\mathrm{FP}}_G$ 和狄拉克算符 $D'$，这两个量都必须用有限个带耦合 $c_i$ 的项来近似和参数化。对于规范作用量，可以使用链变量的闭合圈上的迹及其幂次的线性组合，
\begin{equation}
  S^{\mathrm{FP}}_G[U] = \sum_m\sum_{l\in\mathcal{L}}c_m(l)\,
    \left(\frac{1}{3}\,\mathrm{Re}\,\tr[1-U_l]\right)^m .
\end{equation}
在这个方程中，$\mathcal{L}$ 是格点上闭合圈的集合，$U_l$ 是沿这样的圈 $l \in \mathcal{L}$ 的链变量的有序乘积。耦合 $c_m(l)$ 是作用量的实参数。圈的集合 $\mathcal{L}$ 必须选择得满足格子对称性，例如离散平移、转动。显然，参数化作用量 (9.59) 是 Wilson 规范作用量 (2.49) 的推广，如果我们在 (9.59) 中选取 $\mathcal{L}$ 为所有基本方格的集合，并只限于 $m = 1$ 的幂次，它就退化为该作用量。

狄拉克算符的参数化以类似的方式进行 [30, 31]，也可以看作 Wilson 狄拉克算符 (5.51) 的推广。当我们在第 5.2 节讨论加倍问题时，被迫包含带狄拉克空间单位矩阵的项。在参数化狄拉克算符的拟设中，现在使用了 Clifford 代数的全部 16 个生成元 $\Gamma_k \in \{1, \gamma_\mu, \sigma_{\mu\nu}, \gamma_\mu\gamma_5, \gamma_5\}$。与规范作用量类似，我们允许一整组连接格子上两个格点 $n$、$m$ 的路径 $p \in \mathcal{P}^{(k)}_{n,m}$，并得到 [30, 31]
\begin{equation}
  D[U]_{nm} = \sum_{k=1}^{16}\Gamma_k
    \sum_{p\in\mathcal{P}^{(k)}_{n,m}}c^{(k,p)}_{n,m}[U]\,U_p ,
\end{equation}
其中 $U_p$ 表示沿路径 $p$ 的链变量乘积。同样，路径集合 $\mathcal{P}^{(k)}_{n,m}$ 必须尊重离散格子对称性。此外，路径集合 $\mathcal{P}^{(k)}_{n,m}$ 和系数 $c^{(k,p)}_{n,m}$ 依赖于 Clifford 代数的元素，因为对于费米子，反射和电荷共轭等其他对称性（见第 5.4 节）也涉及作用在狄拉克空间中的矩阵。系数 $c^{(k,p)}_{n,m}[U]$ 可以是链变量 $U_\mu$ 的规范不变函数，例如它们可以依赖于局部基本方格的值。

对于规范作用量 $S^{\mathrm{FP}}_G$，(9.54) 两端的耦合相同，并确定不动点作用量，然后它被用于任意的 $\beta$。实际的确定通过极小化 (9.54) 左端与右端的差来完成。这一步实际上涉及两个极小化：作用量的系数和细位形都必须变化。在一种实用的方法 [26--29, 32] 中，首先在 (9.54) 的右端使用一个较简单的规范作用量，它用于确定与粗位形 $U'$ 对应的极小规范位形 $U^{\min}$。第二步再用这个极小位形确定参数化不动点作用量 (9.59) 的系数 $c_m(l)$。所得的不动点规范作用量经过了非微扰测试，并被证明只有很小的截断效应 [26--29]。

对于费米子 RG 方程 (9.58)，选择了一种与规范场不动点作用量略有不同的方法 [33--37]。不是为 RG 方程寻找不动点解，而是用 (9.58) 确定一个狄拉克算符序列：从一个具有非常大 $\beta$ 和狄拉克算符 $D$ 的细格子出发，其中耦合保持固定。然后由 (9.58) 确定粗格子新狄拉克算符 $D'$ 的耦合。迭代这一步骤，可以确定目标 $\beta$ 处的狄拉克算符，它给出人们想要工作的分辨率。每一步的系数都通过极小化一个可选取为如下形式的 $\chi^2$ 泛函由 (9.58) 计算
\begin{equation}
  \chi^2 = \min\sum_i\sum_k
    \left\|\left(D'[U'_{(i)}] - D_{\mathrm{rhs}}[U^{\min}_{(i)}]\right)
      \psi^{\mathrm{rand}}_{(k)}\right\|^2 .
\end{equation}
在这个方程中，$D'[U'_{(i)}]$ 是粗格子上的狄拉克算符，$D_{\mathrm{rhs}}[U^{\min}_{(i)}]$ 表示 RG 方程 (9.58) 的右端。指标为 $i$ 的求和跑遍一个粗规范位形系综 $U'_{(i)}$（$i = 1, 2, \ldots$），由 FP 方程 (9.54) 确定的相应极小位形记作 $U^{\min}_{(i)}$。两个狄拉克算符的差然后作用在一组由指标 $k = 1, 2, \ldots$ 标记的随机矢量 $\psi^{\mathrm{rand}}_{(k)}$ 上。狄拉克算符差作用在 $\psi^{\mathrm{rand}}_{(k)}$ 上产生一组矢量，$\|\ldots\|$ 表示这些矢量的范数。

所得的 FP 狄拉克算符在一系列论文 [32--39] 中得到了测试，并确认诸如小截断效应等预期性质可以被参数化 (9.59) 和 (9.60) 所描述。也许更重要的是，RG 方程把连续理论的对称性输运到格点上。对于手征对称性，这些性质在 [35, 40, 41] 中得到了详细讨论。特别是，FP 狄拉克算符满足 Ginsparg--Wilson 方程，并且我们在第 7 章已经讨论过的性质，如正确的反常、指标定理，随之而来。我们将在下一节重新讨论 RG 变换把对称性输运到格点这一事实。

我们指出，狄拉克算符的参数化 (9.60) 还被以另一种方式使用 [38, 42]。把参数化狄拉克算符 (9.60) 插入 Ginsparg--Wilson 方程 (7.29)，可以读出系数 $c^{(k,p)}_{n,m}$ 的一组耦合二次方程组。在截断之后，这组方程可以数值求解，所得的系数确定所谓的近手征改进狄拉克算符。这个算符经过了广泛测试 [39, 43]，被发现具有良好的手征性质，而数值开销比重叠算符（overlap operator）所需的低得多。

\section{把连续对称性映射到格点}
\label{sec:9.4}

在上一节中，我们看到，对完整格点 QCD 实际实施 RG 方法是一项相当不平凡的事业。我们以实空间重整化技术的另一个非常重要的应用来结束本章：我们证明阻塞变换可以用来把连续作用量的对称性映射到由阻塞得到的相应格点作用量的对称性的格点版本。这里我们只对费米子情形讨论连续对称性与格点对称性之间的联系。这一情形特别重要，因为所关心的对称性之一是手征对称性——格点表述的一个长期难题。如果把连续理论中的（矢量型）手征对称性映射到格点上，会发现格点上涌现的结构就是我们在第 7 章已详细讨论过的 Ginsparg--Wilson 方程。然而，连续对称性与其格点对应物之间的关系并不限于手征对称性，而是可以对连续理论的一般对称性推导出来。

\subsection{生成泛函及其对称性}
\label{sec:9.4.1}

与第 9.2 节讨论的自由格点费米子的阻塞类似，我们分析的出发点又是那个通过从费米子连续作用量 $S$ 出发的阻塞变换定义格点上狄拉克算符 $D$ 的方程：
\begin{equation}
  e^{-\bar\psi D\psi} = \int D[\varphi,\bar\varphi]\,
    e^{-(\bar\psi-\bar\varphi^B)B(\psi-\varphi^B) - S_F[\varphi,\bar\varphi]} .
\end{equation}
我们用 $\psi$、$\bar\psi$ 表示格点费米子，并对所有指标（时空、色、狄拉克）使用矢量/矩阵记号。$\varphi$、$\bar\varphi$ 是连续理论中进入右端路径积分的费米子场。通过对以格点为中心的超立方体积分，从它们构造出阻塞场 $\varphi^B$、$\bar\varphi^B$。后者生活在格子的格点上，因此它们与格点场 $\psi$、$\bar\psi$ 具有相同的指标，特别是离散的时空指标。

我们强调，(9.62) 是在背景规范场中理解的，该规范场不被积掉，但为了保持规范协变性也必须以适当的方式阻塞（见上一节）。(9.62) 中的连续路径积分只是形式地定义的，但例如可以通过上一节所做的那样由连续极限中的格点理论来构造。无论如何，这里我们只对作用量中费米子部分（它是费米子场的双线性形式）的对称性性质感兴趣。对于这一分析，连续路径积分的形式定义就足够了。

这里我们考虑比第 9.2 节所用更一般的阻塞方案，允许一个非平凡的阻塞核 $B$，它决定阻塞场 $\varphi^B$、$\bar\varphi^B$ 与格点场 $\psi$、$\bar\psi$ 如何混合。目前我们暂不指定阻塞核 $B$，但会在本节末尾讨论它的作用。

Ginsparg 和 Wilson 的工作 [44] 从 (9.62) 出发，分析它在格点场手征转动下的行为。按照 [45]，我们使用一个略有不同的方法，考虑如下定义的格点上的生成泛函
\begin{equation}
  W[\bar J,J] = \int D[\psi,\bar\psi]\,e^{\bar\psi J + \bar J\psi - \bar\psi D\psi} ,
\end{equation}
其中我们有耦合到格点费米子的源 $J$ 和 $\bar J$。插入来自 (9.62) 的作用量指数，我们通过阻塞方案找到一个生成泛函的表达式：
\begin{align}
  W[\bar J,J] &= \int D[\psi,\bar\psi]\,e^{\bar\psi J + \bar J\psi}
    \int D[\varphi,\bar\varphi]\,
    e^{-(\bar\psi-\bar\varphi^B)B(\psi-\varphi^B) - S_F[\varphi,\bar\varphi]} \nonumber\\
  &= \int D[\varphi,\bar\varphi]\,e^{-\bar\varphi^B B\varphi^B - S_F[\varphi,\bar\varphi]}
    \int D[\psi,\bar\psi]\,
    e^{-\bar\psi B\psi + \bar\psi(J+B\varphi^B) + (\bar J+\bar\varphi^B B)\psi} \nonumber\\
  &= \det[B]\,e^{\bar J B^{-1} J}\int D[\varphi,\bar\varphi]\,
    e^{\bar J\varphi^B + \bar\varphi^B J - S_F[\varphi,\bar\varphi]} .
\end{align}
在最后一步中，我们已经通过在指数中配平方（比较第 9.2 节的讨论）形式地求解了对格点场的高斯积分。所得的表达式 (9.64) 是用连续路径积分表示的格点生成泛函。

我们现在探讨连续费米子作用量 $S_F$ 的对称性如何影响格点生成泛函 $W[\bar J,J]$。特别地，我们考虑连续场的如下变换：
\begin{equation}
  \varphi \to \varphi' = e^{i\varepsilon T}\varphi , \qquad
  \bar\varphi \to \bar\varphi' = \bar\varphi\,e^{i\varepsilon T} .
\end{equation}
变换的生成元 $T$、$T$ 是作用在狄拉克空间中的矩阵。更一般的变换是可能的，正如 [46] 中所讨论的。我们的选择意味着阻塞场 $\varphi^B$ 和 $\bar\varphi^B$ 与原始场以相同的方式变换，因为从 $\varphi$、$\bar\varphi$ 到 $\varphi^B$、$\bar\varphi^B$ 的阻塞是纯粹标量操作，换言之，阻塞场本质上是连续场的线性组合。我们强调，生成元 $T$ 和 $\bar T$ 是独立的变换，即 $\varphi$ 和 $\bar\varphi$ 不必以相同的方式变换。

我们现在假定变换 (9.65) 是作用量的对称性
\begin{equation}
  S_F[\varphi',\bar\varphi'] = S_F[\varphi,\bar\varphi] .
\end{equation}
我们也可以使用变换后的变量 $\varphi'$、$\bar\varphi'$ 来计算 (9.64) 最后一行中对连续场的积分。这样做并利用作用量的不变性，我们得到
\begin{align}
  &\int D[\varphi',\bar\varphi']\,e^{\bar J\varphi'^B + \bar\varphi'^B J - S_F[\varphi',\bar\varphi']} \nonumber\\
  &= \int D\left[e^{i\varepsilon T}\varphi,\,\bar\varphi\,e^{i\varepsilon T}\right]\,
     \exp\!\left[\bar J e^{i\varepsilon T}\varphi^B
       + \bar\varphi^B e^{i\varepsilon T}J\right. \nonumber\\
  &\qquad \left.- S_F[\varphi,\bar\varphi]\right] \nonumber\\
  &= \left(1 + i\varepsilon A_{TT} + O(\varepsilon^2)\right)
     \int D[\varphi,\bar\varphi]\,
     \exp\!\left[\bar J e^{i\varepsilon T}\varphi^B
       + \bar\varphi^B e^{i\varepsilon T}J\right. \nonumber\\
  &\qquad \left.- S_F[\varphi,\bar\varphi]\right] .
\end{align}
在最后一步中，我们变换了连续路径积分的测度，并考虑到变换 (9.65) 可能是反常的，反常 $A_{TT}$ 出现在变换的雅可比行列式中 [47, 48]。由于稍后我们将把一切表达式求值到 $\order{\varepsilon}$，我们只保留了雅可比行列式的首阶项。对于无反常的变换 $T$、$\bar T$，有 $A_{TT} = 0$。

把结果 (9.67) 代回生成泛函的表达式 (9.64)，我们发现连续对称性 (9.65) 和 (9.66) 蕴含格点生成泛函的如下对称性：
\begin{equation}
  W[\bar J,J] = e^{\bar J\left(B^{-1} - e^{i\varepsilon T}B^{-1}e^{i\varepsilon T}\right)J}
    \left(1 + i\varepsilon A_{TT} + O(\varepsilon^2)\right)
    W\left[e^{i\varepsilon T}J,\,\bar J e^{i\varepsilon T}\right] .
\end{equation}
方程 (9.68) 总结了连续对称性如何在通过阻塞构造的格点生成泛函 $W[\bar J,J]$ 中体现出来。

\subsection{识别相应的格点对称性}
\label{sec:9.4.2}

在分析了连续对称性在我们格点生成泛函中的体现之后，我们现在想要识别格点上一个相应的、也满足对称性条件 (9.68) 的对称性。为此我们考虑变换后的格点场
\begin{equation}
  \psi' = e^{i\varepsilon M}\psi , \qquad
  \bar\psi' = \bar\psi\,e^{i\varepsilon M} .
\end{equation}
格点变换的两个生成元 $M$、$M$ 尚不知道，我们想要识别它们如何依赖于连续生成元 $T$、$\bar T$ 和阻塞核 $B$。该变换应是格点作用量的对称性这一事实蕴含
\begin{equation}
  \bar\psi' D\psi' = \bar\psi\,e^{i\varepsilon M}D\,e^{i\varepsilon M}\psi = \bar\psi D\psi .
\end{equation}
为便于后文使用，我们指出，当按 $\varepsilon$ 展开时，不变性条件 (9.70) 在 $\order{\varepsilon}$ 阶蕴含对易关系
\begin{equation}
  M D + D M = 0 .
\end{equation}
正如我们在上一节对连续表达式所做的那样，我们现在用变换后的场 $\psi'$、$\bar\psi'$ 表示生成泛函 (9.63)，并探讨对称性 (9.70) 的含义：
\begin{align}
  W[\bar J,J] &= \int D[\psi',\bar\psi']\,e^{-\bar\psi' D\psi' + \bar\psi' J + \bar J\psi'} \nonumber\\
  &= \det\!\left[e^{i\varepsilon M}\right]\det\!\left[e^{i\varepsilon M}\right]
     \int D[\psi,\bar\psi]\,
     e^{-\bar\psi D\psi + \bar\psi e^{i\varepsilon M}J + \bar J e^{i\varepsilon M}\psi} .
\end{align}
右端的两个雅可比行列式来自格点测度的变换。利用关系 $\det[A] = \exp(\tr\ln A)$（见 (A.54)），我们可以把它们展开为
\begin{equation}
  \det\!\left[e^{i\varepsilon M}\right]\det\!\left[e^{i\varepsilon M}\right]
  = 1 + i\varepsilon\,\tr[M + M] + O(\varepsilon^2) .
\end{equation}
结合最后两个方程，我们发现由格点对称性 (9.70) 蕴含的生成泛函对称性关系：
\begin{equation}
  W[\bar J,J] = \left(1 + i\varepsilon\,\tr[M + M] + O(\varepsilon^2)\right)
    W\left[e^{i\varepsilon M}J,\,\bar J e^{i\varepsilon M}\right] .
\end{equation}
现在我们把来自格点变换的对称性条件 (9.74) 与相应的连续关系 (9.68) 进行比较。令两者相等，我们得到方程
\begin{multline}
  e^{\bar J\left(B^{-1} - e^{i\varepsilon T}B^{-1}e^{i\varepsilon T}\right)J}
    \left(1 + i\varepsilon A_{TT} + O(\varepsilon^2)\right)
    W\left[e^{i\varepsilon T}J,\,\bar J e^{i\varepsilon T}\right] \\
  = \left(1 + i\varepsilon\,\tr[M + M] + O(\varepsilon^2)\right)
    W\left[e^{i\varepsilon M}J,\,\bar J e^{i\varepsilon M}\right] ,
\end{multline}
它可以用来识别与连续变换 $T$、$\bar T$ 匹配的格点变换 $M$、$M$。最后一步是插入生成泛函的显式形式，
\begin{equation}
  W[\bar J,J] = \det[D]\,e^{\bar J D^{-1}J} ,
\end{equation}
它通过直接求解高斯积分 (9.63) 得到。当插入 (9.76) 时，(9.75) 变成（我们略去两端的因子 $\det[D]$）
\begin{multline}
  e^{\bar J\left(B^{-1} - e^{i\varepsilon T}B^{-1}e^{i\varepsilon T}\right)J}
    \left(1 + i\varepsilon A_{TT} + O(\varepsilon^2)\right)
    e^{\bar J e^{i\varepsilon T}D^{-1}e^{i\varepsilon T}J} \\
  = \left(1 + i\varepsilon\,\tr[M + M] + O(\varepsilon^2)\right)
    e^{\bar J e^{i\varepsilon M}D^{-1}e^{i\varepsilon M}J} .
\end{multline}
最后一个方程对任意的 $\varepsilon$ 和任意的源 $J$、$\bar J$ 成立。因此 $J$ 和 $\bar J$ 的双线性项，以及与 $J$、$\bar J$ 无关的 $\order{\varepsilon}$ 项必须匹配。对于后一项我们得出结论
\begin{equation}
  A_{TT} = \tr[M + M] ,
\end{equation}
从而识别出连续反常的格点对应物。源 $J$、$\bar J$ 的双线性项导致夸克传播子的一个对称性关系，
\begin{equation}
  B^{-1} - e^{i\varepsilon T}B^{-1}e^{i\varepsilon T} + e^{i\varepsilon T}D^{-1}e^{i\varepsilon T}
  = e^{i\varepsilon M}D^{-1}e^{i\varepsilon M} .
\end{equation}
当按 $\varepsilon$ 展开时，在 $\order{\varepsilon}$ 阶得到
\begin{equation}
  T\left(D^{-1} - B^{-1}\right) + \left(D^{-1} - B^{-1}\right)T
  = M D^{-1} + D^{-1}M .
\end{equation}
这个方程由
\begin{equation}
  M = T\left(1 - B^{-1}D\right) , \qquad
  M = \left(1 - DB^{-1}\right)T .
\end{equation}
解出。这些就是我们想要找到的格点对称性的生成元。它们依赖于连续生成元 $T$、$\bar T$、阻塞核 $B$ 和格点狄拉克算符 $D$。把 $M$、$M$ 插入对称性关系 (9.71)，最终得到格点狄拉克算符 $D$ 的一个非线性方程，
\begin{equation}
  T D + D T = D\left[B^{-1}T + T B^{-1}\right]D ,
\end{equation}
它是 Ginsparg--Wilson 方程 [45, 46] 的推广。

把显式形式 (9.81) 插入反常方程 (9.78)，得到反常的最终形式
\begin{equation}
  A_{TT} = \tr\left[T + T - T B^{-1}D - D B^{-1}T\right] .
\end{equation}
容易检验，当用阻塞核 $B = 2\cdot\mathbb{1}$ 阻塞一个矢量型连续理论，并考虑连续理论中 $T = \bar T = \gamma_5$ 的手征转动时，方程 (9.82) 退化为通常的 Ginsparg--Wilson 关系 (7.29)。(9.81) 的生成元 $M$、$M$ 就是 L\"uscher 对称性的生成元 (7.2)，反常取形式 $A = -\tr[\gamma_5 D]$，这我们在第 7 章已经讨论过。

我们已经证明，借助阻塞变换，我们可以识别在格点上实现手征对称性的数学结构，特别是 Ginsparg--Wilson 方程、格点上的手征变换以及反常的形式。所有这些在第 7 章中都是分别发展的，但现在从单次计算中一并得出。最近，阻塞变换在正确实现格点外尔费米子（标准模型电弱部分所需的）的动机下被进一步探索 [46, 49--51]。一个核心见解 [50] 涉及阻塞矩阵 $B$ 的作用：
\begin{enumerate}
  \item 阻塞矩阵 $B$ 必须破缺目标理论中反常的所有对称性。
  \item 其他对称性如果方便也可以被 $B$ 破缺。
\end{enumerate}
对于 QCD 的情形，我们在上面讨论中使用的 $B \propto \mathbb{1}$ 的选择就足够了，因为只有 U(1) 轴变换是反常的。对于电弱部分的情形，由于费米子数反常，需要额外的破缺。对于这一情形，在某个特定模型中提出了合适的阻塞矩阵 [50]，并证明对于相应的不动点作用量会涌现出正确的对称性模式。对于所得出的 Ginsparg--Wilson 方程，有可能使用重叠构造找到显式解 [51]。

\section*{参考文献}

\begin{enumerate}[label={[\arabic*]}]
  \item K.~Symanzik: Nucl.\ Phys.\ B 226, 187 (1983).
  \item K.~Symanzik: Nucl.\ Phys.\ B 226, 205 (1983).
  \item M.~L\"uscher and P.~Weisz: Commun.\ Math.\ Phys.\ 97, 59 (1985).
  \item M.~L\"uscher and P.~Weisz: Erratum, Commun.\ Math.\ Phys.\ 98, 433 (1985).
  \item B.~Sheikholeslami and R.~Wohlert: Nucl.\ Phys.\ B 259, 572 (1985).
  \item M.~L\"uscher, S.~Sint, R.~Sommer, and P.~Weisz: Nucl.\ Phys.\ B 478, 365 (1996).
  \item M.~L\"uscher: in \emph{Probing the Standard Model of Particle
        Interactions}, Proceedings Les Houches 1997, edited by R.~Gupta,
        A.~Morel, E.~de Rafael, and F.~David, Vol.~1 (Elsevier, Amsterdam
        1998).
  \item G.~Curci, P.~Menotti, and G.~Paffuti: Phys.\ Lett.\ B 130, 205 (1983).
  \item G.~Curci, P.~Menotti, and G.~Paffuti: Erratum, Phys.\ Lett.\ B 135, 516 (1984).
  \item R.~Wohlert (unpublished) DESY 87--069 (1987).
  \item M.~L\"uscher and P.~Weisz: Nucl.\ Phys.\ B 479, 429 (1996).
  \item G.~P.~Lepage and P.~B.~Mackenzie: Phys.\ Rev.\ D 48, 2250 (1993).
  \item T.~DeGrand and C.~DeTar: \emph{Lattice Methods for Quantum
        Chromodynamics} (World Scientific, Singapore 2006).
  \item K.~Symanzik: Nucl.\ Phys.\ B 190, 1 (1981).
  \item M.~L\"uscher: Nucl.\ Phys.\ B 254, 252 (1985).
  \item M.~L\"uscher, R.~Narayanan, P.~Weisz, and U.~Wolff: Nucl.\ Phys.\ B 384, 168 (1992).
  \item S.~Sint: Nucl.\ Phys.\ B 421, 135 (1994).
  \item S.~Sint: Nucl.\ Phys.\ B 451, 416 (1995).
  \item M.~L\"uscher et al.: Nucl.\ Phys.\ B 491, 323 (1997).
  \item R.~G.~Edwards, U.~M.~Heller, and T.~R.~Klassen: Nucl.\ Phys.\ (Proc.\ Suppl.) 63, 847 (1998).
  \item K.~Jansen and R.~Sommer: Nucl.\ Phys.\ B 530, 185 (1998).
  \item W.~Bietenholz and U.-J.~Wiese: Nucl.\ Phys.\ B 464, 319 (1996).
  \item K.~G.~Wilson and J.~Kogut: Phys.\ Rep.\ 12C, 76 (1974).
  \item K.~G.~Wilson: Rev.\ Mod.\ Phys.\ 47, 773 (1975).
  \item P.~Hasenfratz: in \emph{Advanced School on Non-Perturbative Quantum
        Field Physics: Peniscola, Spain 2--6 June 1997}, edited by
        M.~Asorey and A.~Dobado, p.~137 (World Scientific, Singapore 1998).
  \item P.~Hasenfratz and F.~Niedermayer: Nucl.\ Phys.\ B 414, 785 (1994).
  \item T.~DeGrand, A.~Hasenfratz, P.~Hasenfratz, and F.~Niedermayer: Nucl.\ Phys.\ B 454, 587 (1995).
  \item T.~DeGrand, A.~Hasenfratz, P.~Hasenfratz, and F.~Niedermayer: Nucl.\ Phys.\ B 454, 615 (1995).
  \item M.~Blatter and F.~Niedermayer: Nucl.\ Phys.\ B 482, 286 (1996).
  \item P.~Hasenfratz et al.: Int.\ J.\ Mod.\ Phys.\ C12, 691 (2001).
  \item C.~Gattringer: Phys.\ Rev.\ D 63, 114501 (2001).
  \item C.~B.~Lang and T.~K.~Pany: Nucl.\ Phys.\ B 513, 645 (1998).
  \item P.~Hasenfratz et al.: Nucl.\ Phys.\ B (Proc.\ Suppl.) 94, 627 (2001).
  \item P.~Hasenfratz et al.: Nucl.\ Phys.\ B (Proc.\ Suppl.) 106, 799 (2002).
  \item P.~Hasenfratz et al.: Nucl.\ Phys.\ B 643, 280 (2002).
  \item S.~Hauswirth: Ph.D.\ thesis, Universit\"at Bern, Switzerland (2002).
  \item T.~J\"org: Ph.D.\ thesis, Universit\"at Bern, Switzerland (2002).
  \item C.~Gattringer: Nucl.\ Phys.\ B (Proc.\ Suppl.) 119, 122 (2002).
  \item C.~Gattringer et al.\ [BGR (Bern-Graz-Regensburg) collaboration]: Nucl.\ Phys.\ B 677, 3 (2004).
  \item P.~Hasenfratz, V.~Laliena, and F.~Niedermayer: Phys.\ Lett.\ B 427, 125 (1998).
  \item P.~Hasenfratz: Nucl.\ Phys.\ B 525, 401 (1998).
  \item C.~Gattringer, I.~Hip, and C.~B.~Lang: Nucl.\ Phys.\ B 597, 451 (2001).
  \item C.~Gattringer et al.: Phys.\ Rev.\ D 79, 054501 (2009).
  \item P.~H.~Ginsparg and K.~G.~Wilson: Phys.\ Rev.\ D 25, 2649 (1982).
  \item C.~Gattringer and M.~Pak: PoS LATTICE2007, 081 (2007).
  \item P.~Hasenfratz, F.~Niedermayer, and R.~von Allmen: JHEP 0610, 010 (2006).
  \item K.~Fujikawa: Phys.\ Rev.\ Lett.\ 42, 1195 (1979).
  \item K.~Fujikawa: Phys.\ Rev.\ D 21, 2848 (1980).
  \item P.~Hasenfratz and M.~Bissegger: Phys.\ Lett.\ B 613, 57 (2005).
  \item P.~Hasenfratz and R.~von Allmen: JHEP 0802, 079 (2008).
  \item C.~Gattringer and M.~Pak: Nucl.\ Phys.\ B 801, 353 (2008).
\end{enumerate}
